\section{模型假设与符号说明}
\subsection{模型假设}

结合题目条件及模型求解需要，作如下假设：

\begin{enumerate}
    \item 问题一中，小区负载、光伏发电功率预测和外部电网电价在制定购电计划时均已知；问题二至问题四严格按照决策时点的信息边界进行求解，尚未获得的实际负载、光伏发电功率及电价不作为提前决策的输入，仅用于实际运行结果及费用计算。

    \item 微网仅考虑从外部电网购电，不考虑向外部电网售电；负载及储能需求以外的剩余供给不产生额外收益，其弃光量和外部电网购电溢余量按后文统一口径进行统计。

    \item 储能设备仅考虑充放电效率造成的能量损耗，不考虑其他内部损耗。充、放电效率均为 $90\%$，储电量保持在 $1200\sim10800\,\mathrm{kWh}$ 范围内，最大充、放电功率均为 $5000\,\mathrm{kW}$。

    \item 同一时段内储能设备不同时进行充电和放电，且各时段的充放电功率在对应 $10\,\mathrm{min}$ 调度区间内保持不变。

    \item 问题一按照题目要求设置
    $E_0=E_T=6000\,\mathrm{kWh}$；
    问题二至问题四不设置每日储电量复位约束，前一日末储电量作为下一日初始储电量，以保持储能状态跨日连续。

    \item 模型不单独设置日照时长变量，光伏发电时段按照各问题所采用的光伏实际数据或预测数据确定。
\end{enumerate}    




\subsection{符号说明}

为统一表述问题一至问题四的购电与储能调度模型，本文采用如下主要符号。

\begin{center}
\begin{tabular}{p{0.18\textwidth}p{0.46\textwidth}p{0.18\textwidth}}
\hline
符号 & 含义 & 单位 \\
\hline
$d$ & 日期索引 & -- \\
$t$ & 时段索引 & -- \\
$T$ & 总时段数 & -- \\
$\Delta t$ & 单个时段长度 & h \\
$P_{d,t}^{L}$ & 小区负载功率 & kW \\
$P_{d,t}^{PV}$ & 光伏发电功率 & kW \\
$L_{d,t}$ & 小区负载电量 & kWh \\
$R_{d,t}$ & 光伏发电电量 & kWh \\
$c_{d,t}$ & 问题一至三采用的分时电价 & 元/kWh \\
$x_{d,t}$ & 正常外部电网购电量 & kWh \\
$e_{d,t}$ & 紧急购电量 & kWh \\
$u_{d,t}$ & 储能充电量 & kWh \\
$v_{d,t}$ & 储能放电量 & kWh \\
$E_{d,t}$ & 时段末储电量 & kWh \\
$w_{d,t}$ & 未利用的光伏电量（弃光量） & kWh \\
$z_{d,t}$ & 储能充放电状态二元变量 & -- \\
$\delta_d$ & 目标日每个10 min时段的安全冗余电量 & kWh \\
$\hat L_{d,t}$ & 负载预测值 & kWh \\
$\hat R_{d,t}$ & 光伏预测值 & kWh \\
$g_{d,t}^{spill}$ & 外部电网购电溢余量 & kWh \\
$w_{d,t}^{PV}$ & 光伏弃光量 & kWh \\
\hline
\end{tabular}
\end{center}
